\chapter{值、类型和初始化}

\section{问题入口：金额为什么不能随便用 double}

假设要写一个购物车，最直接的金额表示是 \code{double}：

\begin{lstlisting}[language=C++,caption={金额直接用 double 的朴素写法}]
double price = 0.1;
double total = price + price + price;
if (total == 0.3) {
    // ...
}
\end{lstlisting}

这段代码的问题不在 \cpp 独有，而在二进制浮点无法精确表示很多十进制小数。\code{0.1 + 0.1 + 0.1} 未必和 \code{0.3} 完全相等。金额通常应该用最小单位的整数，例如“分”。

\begin{lstlisting}[language=C++,caption={用整数分表示金额}]
#include <cstdint>

struct Money {
    std::int64_t cents = 0;
};

Money add(Money lhs, Money rhs)
{
    return Money{.cents = lhs.cents + rhs.cents};
}
\end{lstlisting}

这里的重点是类型不是装饰。类型表达业务含义，也限制错误。一个普通 \code{int} 可能是年龄、下标、数量、分数或错误码；而 \code{Money} 明确告诉读者它表示金额。

\section{初始化：少给错误留入口}

\cpp 有多种初始化形式。入门阶段建议优先使用大括号初始化，因为它能挡住一部分窄化转换。

\begin{lstlisting}[language=C++,caption={大括号初始化拒绝窄化转换}]
int ok{42};
// int bad{3.14}; // 编译错误：double 到 int 会丢失信息
int truncated = 3.14; // 可以编译，但值变成 3
\end{lstlisting}

不要把这理解成“永远只用一种形式”。构造对象、调用工厂函数、使用标准库容器时会遇到不同写法。先建立原则：初始化应该让对象一出生就处在有效状态。

\section{const 不是形式主义}

\code{const} 的意义是把“不修改”写进接口。下面这个函数只观察价格列表，不应该复制，也不应该修改。

\begin{lstlisting}[language=C++,caption={用 const 引用表达只读访问}]
#include <vector>

std::int64_t total_cents(const std::vector<Money>& items)
{
    std::int64_t total = 0;
    for (const Money& item : items) {
        total += item.cents;
    }
    return total;
}
\end{lstlisting}

如果写成 \code{std::vector<Money> items}，每次调用都会复制整个数组。复制不一定错误，但这里没有必要。接口应该表达意图：观察就用 \code{const&} 或 \code{std::span<const T>}，拥有或需要本地副本时才按值传递。

\section{auto 的正确用法}

\code{auto} 可以减少重复类型，但不能掩盖所有权和复制。

\begin{lstlisting}[language=C++,caption={range-for 中避免不必要复制}]
std::vector<std::string> names = {"Ada", "Bjarne", "Grace"};

for (const auto& name : names) {
    // 只观察，不复制 string
}

for (auto name : names) {
    // 每轮复制一个 string，除非确实需要本地副本
}
\end{lstlisting}

教学阶段要养成一个习惯：看到 \code{auto}，仍然要问它推导出的到底是值、引用还是指针。语法短不等于语义简单。

\section{边界：整数溢出和符号混用}

整数类型也有坑。\code{int} 通常是 32 位，累加大数据时可能溢出。\code{std::size_t} 是无符号类型，和 \code{int} 混用时可能产生意外比较。

\begin{lstlisting}[language=C++,caption={避免有符号和无符号混乱}]
std::vector<int> values{1, 2, 3};
for (std::size_t i = 0; i < values.size(); ++i) {
    // i 和 size() 类型一致
}
\end{lstlisting}

但如果下标要向左移动到 \code{-1} 表示结束，\code{std::size_t} 就不适合。类型选择不是背规则，而是看值域和操作。

\begin{warningbox}
类型错误常常不是编译错误，而是建模错误。金额、长度、数量、索引、标识符和错误码不要长期混在同一种裸整数里。
\end{warningbox}

\section{从金额例子继续推到类型设计}

只把金额改成整数分还不够。下面这段代码仍然可能把“分”和“件数”混起来：

\begin{lstlisting}[language=C++,caption={裸整数仍然容易混淆含义}]
std::int64_t cents = 1299;
int quantity = 3;
std::int64_t wrong = cents + quantity; // 编译器不知道业务上不合理
\end{lstlisting}

把金额包装成 \code{Money} 后，错误会更难发生。再进一步，可以把常见操作写成函数：

\begin{lstlisting}[language=C++,caption={让金额操作表达业务含义}]
struct Money {
    std::int64_t cents = 0;
};

Money multiply(Money price, int quantity)
{
    return Money{.cents = price.cents * quantity};
}

Money add(Money lhs, Money rhs)
{
    return Money{.cents = lhs.cents + rhs.cents};
}
\end{lstlisting}

这不是为了让代码变繁琐，而是把业务边界交给类型系统。以后如果要禁止负金额、处理货币种类、做舍入策略，修改点会集中在 \code{Money} 附近，而不是散落在所有 \code{int} 计算里。

\section{值、引用和视图的选择}

函数参数最常见的三种形态是按值、按引用和视图。可以用一个用户列表作为判断样本：

\begin{lstlisting}[language=C++,caption={三种参数形态表达不同意图}]
struct User {
    int id = 0;
    std::string name;
};

void rename(User& user, std::string new_name)
{
    user.name = std::move(new_name);
}

void print_user(const User& user)
{
    std::cout << user.id << " " << user.name << '\n';
}

std::size_t total_name_length(std::span<const User> users)
{
    std::size_t total = 0;
    for (const User& user : users) {
        total += user.name.size();
    }
    return total;
}
\end{lstlisting}

\code{rename} 要修改已有对象，所以用非 \code{const} 引用。\code{print_user} 只观察一个对象，所以用 \code{const&}。\code{total_name_length} 观察一段连续对象，所以用 \code{std::span<const User>}。这三个接口的区别不是语法偏好，而是所有权和修改权限不同。

\section{初始化和有效状态}

初始化的目标是让对象一出生就可用。反例是先创建半成品，再靠后续赋值补齐：

\begin{lstlisting}[language=C++,caption={半成品对象容易被提前使用}]
Money price;
// 中间几十行代码
price.cents = 1299;
\end{lstlisting}

更好的写法是直接给出有效值：

\begin{lstlisting}[language=C++,caption={直接构造有效值}]
Money price{.cents = 1299};
\end{lstlisting}

如果值需要校验，就不要暴露裸字段，而是放到类的构造函数或工厂函数里。也就是说，初始化不是“把值塞进去”，而是“建立一个满足规则的对象”。

\begin{exercisebox}
设计一个 \code{Percent} 类型表示折扣百分比，要求范围是 \code{0} 到 \code{100}。先写一个能被误用的裸 \code{int} 版本，再把它改成类型。重点观察：哪些错误从运行期判断提前到了构造边界。
\end{exercisebox}
